\scanpage{182}
Pratiq$^{\mathrm{mt}}$, \unsure{cel.|cela|on} revient à contracter $i_0(k_0),i_1(k_1),i_2(k_2)$ en des pts, et \unsure{de} se borner à des \unsure{cochaines|cochaînes} normalisées sur \deleted{\unsure{le complexe}} \inserted{l'\unsure{espace|ensemble} simplicial} \uline{ponctué} obtenu
\diagram{\begin{tikzcd}[ampersand replacement=\&,column sep=1em]
\overline K_0 \& \overline K_1\arrow[l,shift left=.6ex]\arrow[l,shift right=.6ex] \& \overline K_2\arrow[l,shift left=1.1ex]\arrow[l]\arrow[l,shift right=1.1ex]
\end{tikzcd}}

On utilise le fait que
\[
\begin{aligned}
\pi_0(k_*)&=e\\
\text{``}\,\pi_1(k_*)&=e\,\text{''}\qquad\text{i.e. }\pi^1(k_*,G)=e
\end{aligned}
\]
\uline{donc tte donnée de descente admet une ponctuation}

et on utilisera le fait
\[
\pi_0(K_*)=e
\]
pour le fait que \deleted{deux ponctuations} tout \uline{automorphisme} d'une structure de descente \uline{ponctuée} est trivial. \deleted{[et \unsure{tt autom} d'une structure de descente, \unsure{pas néc.} ponctuée, est \unsure{définie|défini} par un élémt de $G$, \unsure{évidemt} unique]}

\deleted{Regardons les}

% archive sheet number: 181
